\section{Stagewise Orthogonality}
\label{app:orthogonality}

The normalized cell has the explicit matrix
\begin{equation}
 T_\theta=
 \begin{bmatrix}
 1&0&c&-s\\
 0&1&s&c\\
 1&0&-c&s\\
 0&-1&s&c
 \end{bmatrix}.
 \label{eq:cell-matrix}
\end{equation}
Its columns have squared norm two.  The first and third columns are orthogonal
because their $c$ contributions cancel; the second and fourth behave
similarly.  The remaining inner products vanish by either sign cancellation or
$cs-sc=0$.  Hence $T_\theta^{\trans}T_\theta=2I$ entrywise.

Direct sums of equal-rank cells retain this Gram identity, and permutations
retain it by conjugation.  The complete Bruun tree also contains shell carry
coordinates because its binomial spine is unbalanced; those coordinates and
the terminal $\sqrt2$ pair weights are included in Definition
\ref{eq:BN-definition}.  The schedule-independent global statement is therefore
obtained from the exact identity $P_NB_N=\DFT_N^{\R}$ and Fourier Parseval:
\begin{equation}
 B_N^{\trans}B_N
 =B_N^{\trans}P_N^{\trans}P_NB_N=NI.
\end{equation}
This argument covers every exact schedule of the same boundary map without
assigning an artificial common depth to the terminal factors.
